\section{Real time formalism for fermions}\label{ss:RTFermion}
In this Appendix, we extend the real time formalism, developed for free bosonic systems in Sec. \ref{ss:RTapproach}, to free fermionic systems.


\subsection{Fermionic system on lattice}
We consider a quadratic Hamiltonian of fermions on lattice labeled by an integer $i=1,\cdots, N$,
\begin{align}
	H = \sum_{i,j = 1}^N \psi^\dagger_i\, M_{ij}\, \psi_j \ ,
\end{align}
with fermionic operators satisfying the anticommutation relation $\{\psi_i, \psi_j^\dagger \} = \delta_{ij}$.
The Hermitian matrix $M$ can be diagonalized to $M=U^\dagger \Delta \,U$ by a unitary matrix $U$, and then the Hamiltonian becomes
\begin{align}
	H = \sum_{l=1}^N\, \lambda_l \,\tilde\psi_l^\dagger\, \tilde\psi_l \ , \qquad \tilde\psi = U\, \psi \ ,
\end{align}
where $\lambda_l$ are real diagonal entries of $\Delta$ and the new operators $\tilde\psi_l$ satisfy the anti-commutation relation $\{\tilde\psi_l, \tilde\psi_m^\dagger \} = \delta_{lm}$.
We let $\CS^{(+)}$ and $\CS^{(-)}$ be the sets of indices for positive and negative eigenvalues, respectively.
We define the ground state as the Dirac sea 
\begin{align}
	|0\rangle \equiv \prod_{l\in \CS^{(-)}} \tilde\psi_l |\tilde 0 \rangle \ ,
\end{align}
where $|\tilde 0 \rangle$ is the state annihilated by all $\tilde\psi_l$.
Then the two-point function of the original fermionic operators $\langle \psi_i \psi_j^\dagger \rangle = \langle 0 | \psi_i \psi_j^\dagger |0 \rangle$ is
\begin{align}\label{FermC}
	C_{ij} \equiv \langle \psi_i \psi_j^\dagger \rangle = \left[ U\, \Theta (\Delta)\, U^\dagger\right]_{ij} \ ,
\end{align}
where $\Theta$ is the step function defined by
\begin{align}
	\Theta (\Delta)_{lm} = \left\{\begin{array}{cl}
						\delta_{lm} & \quad \text{for}\quad l, m \in \CS^{(+)} \ , \\
						0 & \quad \text{otherwise} \ .
					\end{array}\right.
\end{align}

Now we consider the entanglement entropy between the subregion $A$ and its compliment.
As in the bosonic case, we assume the reduced density matrix $\rho_A$ is generated by the modular Hamiltonian $H_A= \sum_{i,j \in A} \psi_i^\dagger\, G_{ij}\, \psi_j$ \cite{2003JPhA...36L.205P} as 
\begin{align}
	\rho_A = \CN\, e^{-H_A} \ ,
\end{align}
where $\CN$ is a normalization constant such that $\tr_A (\rho_A) = 1$.
We diagonalize $G$ by a unitary matrix $V$ and rewrite the modular Hamiltonian as $H_A= \sum_{I} \epsilon_I \, d_I^\dagger d_I$ where new operators $d_I = V_{Ii} \,\psi_i$ for $i\in A$ satisfy $\{ d_I, d_J^\dagger \} = \delta_{IJ}$.
In this basis, the normalization constant $\CN$ is fixed to be $\CN = \prod_{I} (1+ e^{-\epsilon_I})^{-1}$.
The two-point function $C_{ij}$ for $i,j\in A$ is evaluated with the modular Hamiltonian, which is diagonalized to be
\begin{align}
	C_{ij} = N\,\tr_A \left( \psi_i \psi_j^\dagger \, e^{-\CH}\right) = V^\dagger_{iI} \,\frac{\delta_{IJ}}{1 + e^{-\epsilon_I}}\, V_{Jj} \ .
\end{align}
Thus the eigenvalues $\epsilon_I$ can be determined by those $\nu_I$ of $C$,
\begin{align}
	\nu_I =\frac{1}{1 + e^{-\epsilon_I}} \ .
\end{align}
The trace of the $n$th power of $\rho_A$ is
\begin{align}
\begin{aligned}
	\tr_A(\rho_A^n) &= \prod_I \frac{1+ e^{-n \epsilon_I}}{(1+ e^{-\epsilon_I})^n}\ , \\
		&= \prod_I \left[ \nu_I^n + (1 - \nu_I)^n \right] \ ,
\end{aligned}
\end{align}
and we obtain the entanglement entropy Eq.\,\eqref{RenyiEE} given by the matrix $C$,
\begin{align}\label{FermCformula}
	S_A = - \tr_A \left[ C \log C + (1-C) \log (1-C) \right] \ .
\end{align}


\subsection{Free massive fermionic fields}
We proceed to deal with free massive fermions whose action is
\begin{align}
	I = \int \d^dx\sqrt{-g} \, \Psi^\dagger \Gamma^0(\i\,  \Gamma^A e_A^M \nabla_M - \i\, m ) \Psi  \ ,
\end{align}
where $A$ is the local Lorentz index in the range $A = 0,\cdots, d-1$ and $M$ stands for the coordinate indices.
The gamma matrices obey the anti-commutation relations $\{ \Gamma_A, \Gamma_B\} = 2\eta_{AB}$ for the metric $\eta = \text{diag} (-1, 1, 1, \cdots, 1)$.
$e^A$ are the vielbein one-form and $\nabla_\mu$ is the spinor covariant derivative defined by $\nabla_M \equiv \partial_M + \omega_M^{AB}[\Gamma_A, \Gamma_B]/8$ with the spin connection two-form $\omega^{AB}$ satisfying $\d e^A + \omega^{A}_{~B}\wedge e^B = 0$.

We concentrate on the case in $d\ge 3$ dimensions. 
See, for example, \textcite{Herzog:2013py} for the implementation in two dimensions.
We set the theory on the radial coordinates Eq.\,\eqref{RadialLattice} and choose the vielbein as
\begin{align}
\begin{aligned}
	e^0 &= \d t \ , \qquad e^{d-1} = \d r \ , \\
	e^a &= r\, \hat e^a\quad (a=1,\cdots,d-2) \ ,
\end{aligned}
\end{align}
where $\hat e^a$ are the vielbein on $\BS^{d-2}$.
The Hamiltonian becomes
\begin{align}
\begin{aligned}
	H = &\int \d\Omega_{d-2}\,\d r\,r^{d-2} \\
		&\cdot \Psi^\dagger \Gamma^0\left[ - \i\,  \Gamma^{d-1} \left(\partial_r + \frac{d-2}{2r}\right) -  \frac{\i}{r} \Gamma^a \hat e_a^\mu \hat \nabla_\mu  - \i\, m\right] \Psi  \ ,
\end{aligned}
\end{align}
where $\hat \nabla_\mu$ and $m$ are the spinor covariant derivative and the coordinate index on $\BS^{d-2}$, respectively.
We want to reduce this to a number of (1+1)-dimensional fermions on $r$.
We choose the basis of the gamma matrices as follows:
\begin{align}
	\Gamma^0 = \i\,{\bf 1} \otimes \sigma_3\ , \qquad 
	\Gamma^{d-1} = {\bf 1} \otimes  \sigma_2 \ , \qquad 		
	\Gamma^a = \gamma^a \otimes \sigma_1\ ,
\end{align}
where $\gamma^a$ and ${\bf 1}$ are the gamma matrices and the identity matrix in $(d-2)$-dimensions satisfying $\{\gamma_a, \gamma_b\} = 2\delta_{ab}$, and $\sigma_i~(i=1,2,3)$ the Pauli matrices.
Then there appears the Dirac operator $\nslash_{\BS^{d-2}}\equiv \gamma^a \hat e_a^\mu \hat \nabla_\mu$ on $\BS^{d-2}$ in the second term in the angle bracket. 
It is known that the eigenfunction $\psi_l$ of the Dirac operator $\nslash_{\BS^{d}}$ on a unit sphere $\BS^{d}$ is labeled by a non-negative integer $l\ge 0$ with eigenvalues \cite{Camporesi:1995fb},
\begin{align}
	\i\,  \nslash_{\BS^{d}} \chi_{l,d}^{(s)} = s \left( l + \frac{d}{2}\right) \chi_{l,d}^{(s)} \ ,
\end{align}
where $s=\pm$ and their degeneracy $g_f(d,l)$ is 
\begin{align}\label{FermiDeg}
	g_f(d,l) = \frac{2^{[d/2]}\, \Gamma(l + d )}{l! \,\Gamma (d)}  \ .
\end{align}
We expand the Dirac fermion by the eigenfunctions
\begin{align}
	\Psi_l^{(s)} =  \chi^{(s)}_{l,d-2} \otimes \left[ r^{1-d/2}\,\psi_l^{(s)}(r)\right] \ ,
\end{align}
so that $\psi_l^{(s)}(r)$ satisfies the anticommutation relation
\begin{align}
	\left\{ \psi_l^{(s)}(r), \left(\psi_{l'}^{(s')} (r')\right)^\dagger\right\} = \i\, \delta_{ss'}\delta_{ll'}\delta (r-r')\ ,
\end{align}
and rewrite the Hamiltonian as the sum over the angular modes,
\begin{align}
	H &= \sum_{l=0}^\infty  \sum_{s=\pm} g_f(d-2,l)\, H_{l}^{(s)} \ , 
\end{align}
with
\begin{align}
\begin{aligned}
	H_{l}^{(s)} =   \int \d r\, \Bigg[ &-\frac{\i}{2} \left(\psi_l^{(s)}(r)\right)^\dagger \sigma_1\, \partial_r \psi_l^{(s)}(r) \\
		&\quad + \frac{\i}{2} \partial_r \left(\psi_l^{(s)}(r)\right)^\dagger \sigma_1\,  \psi_l^{(s)}(r) \\
	&\qquad +  \frac{s}{r} \left( l + \frac{d-2}{2}\right) \left(\psi_l^{(s)}(r)\right)^\dagger  \sigma_2\, \psi_l^{(s)}(r) \\
	&\qquad\quad + m\, \left(\psi_l^{(s)}(r)\right)^\dagger \sigma_3\,  \psi_l^{(s)}(r) \Bigg]   \ .
\end{aligned}
\end{align}
To discretize the radial coordinate $r$ to lattice labeled by $j=1,2,\cdots, N$ with lattice spacing $a$, we replace 
\begin{align}
	\begin{aligned}
		r &\to j a \ , \qquad &  \delta (r - r') &\to \frac{\delta_{jk}}{a} \ , \\
		\psi_l^{(s)}(r) &\to \frac{ \psi^{(s)}_{l,j} }{\sqrt{a}}\ , \qquad & \partial_r\, \psi_l^{(s)} (r) &\to \frac{\psi^{(s)}_{l,j+1} - \psi^{(s)}_{l,j}}{a}  \ , 
	\end{aligned}
\end{align}
and then the discretized Hamiltonian reads
\begin{align}
	H_l^{(s)} = \frac{1}{a}\sum_{j,k=1}^N \left(\psi^{(s)}_{l,j}\right)^\dagger \, \left(M_l^{(s)}\right)^{j,k} \,\psi^{(s)}_{l,k} 
\end{align}
with the Hermitian matrix $M_l^{(s)}$ given by
\begin{align}\label{HermFermiM}
	\begin{aligned}
		\left(M_l^{(s)}\right)^{k,k} &= \frac{s}{k} \left( l + \frac{d-2}{2}\right) \sigma_2  + (ma)\sigma_3 \ , \\
		\left(M_l^{(s)}\right)^{k,k+1} &= - \left(M_l^{(s)}\right)^{k+1, k} = - \frac{\i}{2}\sigma_1 \ .
	\end{aligned}
\end{align}

Finally the entanglement entropy of the free massive Dirac fermion is given by the sum over the angular modes $l$,
\begin{align}
	S_A = \sum_{l=0}^\infty  \sum_{s=\pm} g_f(d-2,l)\, S_l^{(s)} \ ,
\end{align}
where $S_l^{(s)}$ is the entropy of the $l^{th}$ mode given by the formula \eqref{FermCformula} with the two-point function $C_l^{(s)}$ of the form \eqref{FermC} for the Hermitian matrix $M_l^{(s)}$, Eq.\,\eqref{HermFermiM}.

\endinput